\begin{textsample}{POS Dim 5 – global\_north\_2024\_07 – Score 13.00 – global\_north\_2024\_07/110870726.txt}  \label{ex:f5_pos_308}
\textbf{Advisory} opinion of The Hague-based \textbf{International} \textbf{Court} of Justice is not binding ; Netanyahu slams \textbf{ruling} as a ' decision of lies ' Presiding \textbf{Judge} Nawaf Salam reads the \textbf{ruling} in the \textbf{International} \textbf{Court} of Justice, or World \textbf{Court} in The Hague, Netherlands on July 19, 2024, as the United Nations top \textbf{court} is delivered a \textbf{non-binding} \textbf{advisory} opinion on the legality of Israel ' s 57-year occupation of lands sought for a Palestinian state.
Photo Credit : AP The U.N. ' s top \textbf{court} on July 19 said Israel ' s decades-long occupation of Palestinian territory was"illegal"and needed to end as soon as possible.
The \textbf{advisory} opinion of The Hague-based \textbf{International} \textbf{Court} of Justice is not binding, but it comes amid mounting concern over Israel ' s war against Hamas sparked by the group ' s brutal October 7 attacks.
Related Stories"The \textbf{court} has found that Israel ' s continued presence in the Palestinian Territories is illegal,"ICJ presiding \textbf{judge} Nawaf Salam said, adding :"Israel must end the occupation as rapidly as possible."The @ @ @ @ @ @ @ @ @ @ cease immediately all new settlement activities and to evacuate all settlers"from occupied land.
Israel ' s policies and practices including the building of new settlements and Israel ' s continued maintainance of a wall between the territories"amount to annexation of large parts"of the occupied territory, it said.
Netanyahu slams \textbf{ruling} as a ' decision of lies ' Israeli Prime Minister Benjamin Netanyahu said the \textbf{court} had made a"decision of lies"."The Jewish people are not occupiers in their own land -- not in our eternal capital Jerusalem, nor in our ancestral heritage of Judea and Samaria ( the occupied West Bank ),"Mr.
Netanyahu said in a statement."No decision of lies in The Hague will distort this historical truth, and similarly, the legality of Israeli settlements in all parts of our homeland can not be disputed."' Extreme danger ' The U.N ' s General Assembly asked the ICJ in late 2022 to give an"\textbf{advisory} opinion"on the"@ @ @ @ @ @ @ @ @ @ in the Occupied Palestinian Territory, including East Jerusalem".
The ICJ held a week-long session in February to hear submissions from countries following the request -- supported by most countries within the Assembly.
During the hearings, most speakers called on Israel to end its 57-year occupation.
They warned a prolonged occupation posed an"extreme danger"to stability in the Middle East and beyond.
But the United States said Israel should not be legally obliged to withdraw without taking its"very real security needs"into account.
Israel did not take part in the oral hearings.
Instead, it \textbf{submitted} a written contribution in which it described the questions the \textbf{court} had been asked as"prejudicial"and"tendentious". ' Ongoing \textbf{violation} ' The General Assembly asked the ICJ to consider two questions.
Firstly, the \textbf{court} should examine the \textbf{legal} consequences of what the U.N. called"the ongoing \textbf{violation} by Israel of the right of the Palestinian people to self-determination".
This @ @ @ @ @ @ @ @ @ @ of the Palestinian territory occupied since 1967"and"\textbf{measures} aimed at altering the demographic composition, character and status of the Holy City of Jerusalem".
In June 1967, Israel crushed some of its Arab neighbours in a six-day war, seizing the West Bank including east Jerusalem from Jordan, the Golan Heights from Syria, and the Gaza Strip and Sinai Peninsula from Egypt.
Israel then began to settle the 70,000 sq. km of seized Arab territory.
The U.N. later declared the occupation of Palestinian territory illegal, and Cairo regained Sinai under its 1979 peace deal with Israel. \textbf{Non-binding} opinion The ICJ was also asked to look into the consequences of what it described as Israel ' s"adoption of related discriminatory legislation and \textbf{measures}".
Secondly, the ICJ was asked to advise on how Israel ' s actions"affect the \textbf{legal} status of the occupation"and what are the consequences for the U.N. and other countries.
The ICJ rules in disputes between states.
Normally @ @ @ @ @ @ @ @ @ @ to \textbf{enforce} them.
In this \textbf{case} however, the opinion is \textbf{non-binding}, although most \textbf{advisory} opinions are in fact \textbf{acted} upon.
The ICJ has previously issued \textbf{advisory} opinions on the legality of Kosovo ' s 2008 declaration of independence from Serbia and apartheid South Africa ' s occupation of Namibia.
It also handed down an opinion in 2004 declaring that parts of the wall erected by Israel in the occupied Palestinian territory were illegal and should be torn down.
Israel has not complied with the ICJ \textbf{ruling}, and has resisted calls to route the barrier along the Green Line, the 1949 Armistice Line established after the end of fighting that accompanied Israel ' s establishment a year earlier.
Comments Comments have to be in English, and in full sentences.
They can not be abusive or personal.
Please abide by our community guidelines for posting your comments.
We have migrated to a new commenting platform.
If you are already a registered user of The Hindu and logged in, you may continue to engage @ @ @ @ @ @ @ @ @ @ account please register and login to post comments.
Users can access their older comments by logging into their accounts on Vuukle.

% matched lemmas: act, advisory, case, court, enforce, international, judge, legal, measure, non-binding, ruling, submit, violation
\end{textsample}
